\documentclass[12pt]{article}
\usepackage{graphicx}
\usepackage{amsmath}
\usepackage{amsfonts}
\usepackage{multirow}
\usepackage[margin=0.6in]{geometry}  

\begin{document}
The attention mechanism can be plainly written as:
\begin{gather*} 
\operatorname{Attention}(Q, K, V)=\operatorname{softmax}\!\left(\frac{QK^\top}{\sqrt{d}}\right) V\\
\end{gather*}

The cross-attention mechanism can be written as:
\begin{align*} 
\operatorname{Cross}&\operatorname{Attention}(X_q, X_c; W_q, W_k, W_v) \\
&= \operatorname{softmax}\!\left(\frac{(X_qW_q)(X_cW_k)^\top}{\sqrt{d}}\right)(X_cW_v)
\end{align*}

The relational-attention mechanism can be written as:
\begin{align*} 
\operatorname{Relational}&\operatorname{Attention}(X_q, X_c; W_q, W_k, V) \\
&= \operatorname{softmax}\!\left(\frac{(X_qW_q)(X_cW_k)^\top}{\sqrt{d}}\right)V
\end{align*}

The adaptable-attention mechanism can be written as:
\begin{align*} 
\operatorname{Adaptable}&\operatorname{Attention}(X_q, X_c; \gamma, W_q, W_k, W_v) \\
&= \operatorname{softmax}\!\left(\frac{(X_qW_q)(X_cW_k)^\top}{\sqrt{d}}\right)\left[(1-\gamma)X_c+\gamma I\right]W_v
\end{align*}

The pairwise transitive inference task can be described by:
\begin{table}[h]
\centering
\begin{tabular}{c c c | c}
\hline
Obs.? & $x_i$ & $x_j$ & $x_i > x_j$ \\
\hline
\multirow{4}{*}{Yes} 
  & A & B & 1 \\
  & B & C & 1 \\
  & B & A & 0 \\
  & C & B & 0 \\
\hline
\multirow{2}{*}{No} 
  & A & C & 1 \\
  & C & A & 0 \\
\hline
\end{tabular}
\caption{Truth table for transitive inference.}
\end{table}

\end{document}